\documentclass[../../main.tex]{subfiles}
\graphicspath{{../../image/}{../../figure/}}
\begin{document}

\section{分界线、不变量与 Lax 公式}\label{section:ch11-分界线不变量与Lax公式}

这里将运用广义特征线理论, 建立\eqref{align:ch11-11-1-1}可容许解的有趣而基本的性质。出发点是一个简单但(如我们将看到)非常有用的恒等式。

考虑两个可容许解 $u$ 与 $u^*$, 对应初值 $u_0$ 与 $u_0^*$。追踪与 $u$ 关联的一条极值后向特征线 $\xi(\cdot)$, 以及与 $u^*$ 关联的一条极值后向特征线 $\xi^*(\cdot)$, 二者从任意固定点 $(x,t)\in(-\infty,\infty)\times(0,\infty)$ 出发。于是 $\xi(\cdot)$ 与 $\xi^*(\cdot)$ 是直线, 沿 $\xi(\cdot)$ $u$ 为常数, 等于 $u(x-,t)$ 或 $u(x+,t)$; 沿 $\xi^*(\cdot)$ $u^*$ 为常数, 等于 $u^*(x-,t)$ 或 $u^*(x+,t)$。特别地, $\dot{\xi}(\tau)=f'(u(x\pm,t))$, $\dot{\xi}^*(\tau)=f'(u^*(x\pm,t))$, $0<\tau<t$。

先对 $u$ 再对 $u^*$ 写出\eqref{align:ch11-11-1-1}, 合并所得两方程, 在顶点为 $(x,t),(\xi(0),0),(\xi^*(0),0)$ 的三角形上积分并应用 Green 公式, 得到恒等式
\begin{align}
&\int_0^t\bigl\{f(u(x\pm,t))-f(u^*(\xi(\tau)-,\tau))-f'(u(x\pm,t))[u(x\pm,t)-u^*(\xi(\tau)-,\tau)]\bigr\}\,\mathrm{d}\tau\nonumber\\
&\quad+\int_0^t\bigl\{f(u^*(x\pm,t))-f(u(\xi^*(\tau)-,\tau))-f'(u^*(x\pm,t))[u^*(x\pm,t)-u(\xi^*(\tau)-,\tau)]\bigr\}\,\mathrm{d}\tau\nonumber\\
&\qquad=\int_{\xi^*(0)}^{\xi(0)}[u_0(y)-u_0^*(y)]\,\mathrm{d}y.\label{align:ch11-11-4-1}
\end{align}
由 $f$ 的凸性, \eqref{align:ch11-11-4-1}左端两个积分均非正。因而 \eqref{align:ch11-11-4-1} 给出两个解的比较, 在后文将有大量应用。

作为\eqref{align:ch11-11-4-1}的第一个应用, 我们用它定位与可容许解 $u$ 关联的分界线。分界线概念由\cref{chapter:ch10-广义特征线}引入。在真正非线性标量守恒律的框架下, 按照\cref{chapter:ch10-广义特征线}的讨论, 分界线无激波, 因而由\cref{theorem:ch11-11-1}, 是 $u$ 为常数的直线。


\begin{theorem}\label{theorem:ch11-11-4-1}
与可容许解 $u$ 关联、初值为 $u_0$、沿其上 $u$ 为常数 $\bar u$ 的分界线, 从 $x$ 轴上的点 $(\bar x,0)$ 出发当且仅当
\begin{align}
\int_{\bar x}^{z}[u_0(y)-\bar u]\,\mathrm{d}y\geqslant0,\quad -\infty<z<\infty.\label{align:ch11-11-4-2}
\end{align}
\end{theorem}

\begin{proof}
假设\eqref{align:ch11-11-4-2}成立。在\eqref{align:ch11-11-4-1}中取 $u^*\equiv\bar u$, $t\in(0,\infty)$, $x=\bar x+tf'(\bar u)$。特别地, $\xi^*(\tau)=\bar x+\tau f'(\bar u)$ 且 $\xi^*(0)=\bar x$。于是由\eqref{align:ch11-11-4-2}, \eqref{align:ch11-11-4-1}右端非负。但此时左端两个积分必均为零, 故 $u(x\pm,t)=\bar u$。于是我们证明了直线 $x=\bar x+tf'(\bar u)$ 在 $[0,\infty)$ 上是无激波特征线, 即与 $u$ 关联的分界线。

反之, 假设直线 $x=\bar x+tf'(\bar u)$ 是与 $u$ 关联的分界线。任取 $z\in(-\infty,\infty)$, 固定 $\tilde u$: 若 $z>\bar x$ 则 $\tilde u<\bar u$, 若 $z<\bar x$ 则 $\tilde u>\bar u$。直线 $z+tf'(\tilde u)$ 与 $\bar x+tf'(\bar u)$ 将在某点 $(x,t)$ 相交, $t>0$。在\eqref{align:ch11-11-4-1}中取 $u^*\equiv\tilde u$, 此时 $\xi(0)=\bar x$, $\xi^*(0)=z$。左端非正, 故
\begin{align}
\int_z^{\bar x}[u_0(y)-\tilde u]\,\mathrm{d}y\leqslant0.\label{align:ch11-11-4-3}
\end{align}
令 $\tilde u\to\bar u$, 即得\eqref{align:ch11-11-4-2}。证毕。
\end{proof}

上述命题对解的重要时间不变量的存在性有推论。

\begin{theorem}\label{theorem:ch11-11-4-2}
假设 $u_0$ 在 $(-\infty,\infty)$ 上可积, 且极大值
\begin{align}
\max_x\int_x^{-\infty}u_0(y)\,\mathrm{d}y=q_-,\qquad \max_x\int_x^\infty u_0(y)\,\mathrm{d}y=q_+\label{align:ch11-11-4-4}
\end{align}
存在。若 $u$ 是初值为 $u_0$ 的可容许解, 则对任意 $t>0$,
\begin{align}
\max_x\int_x^{-\infty}u(y,t)\,\mathrm{d}y=q_-,\qquad \max_x\int_x^\infty u(y,t)\,\mathrm{d}y=q_+.\label{align:ch11-11-4-5}
\end{align}
\end{theorem}

\begin{proof}
注意 $q_-$ 存在当且仅当 $q_+$ 存在。事实上, 由\cref{theorem:ch11-11-4-1}, 两个极大均在使直线 $x=\bar x+tf'(0)$ 成为与 $u$ 关联、沿其上 $u$ 恒为零的分界线的 $\bar x$ 集上达到。于是再次由\cref{theorem:ch11-11-4-1}, \eqref{align:ch11-11-4-5}中两个极大均在 $\hat x=\bar x+tf'(0)$ 处达到。

现在归一化 $f$ 使 $f(0)=0$, 先在区域 $\{(y,\tau):0<\tau<t,\ -\infty<y<\bar x+\tau f'(0)\}$ 上、再在区域 $\{(y,\tau):0<\tau<t,\ \bar x+\tau f'(0)<y<\infty\}$ 上积分\eqref{align:ch11-11-1-1}。应用 Green 公式, 并注意 $u$ 沿直线 $x=\bar x+\tau f'(0)$ 为零, 得
\begin{align}
\int_{\hat x}^{-\infty}u(y,t)\,\mathrm{d}y=\int_{\bar x}^{-\infty}u_0(y)\,\mathrm{d}y,\qquad
\int_{\hat x}^\infty u(y,t)\,\mathrm{d}y=\int_{\bar x}^\infty u_0(y)\,\mathrm{d}y,\label{align:ch11-11-4-6}
\end{align}
这验证了\eqref{align:ch11-11-4-5}。证毕。
\end{proof}

真正非线性标量守恒律最引人注目的特征之一, 是可容许解可由初值通过如下程序显式确定。我们从 Legendre 变换开始
\begin{align}
g(v)=\max_u[uv-f(u)],\label{align:ch11-11-4-7}
\end{align}
注意最大值在 $u=[f']^{-1}(v)$ 处达到。对给定初值 $u_0(\cdot)$, 关联\textbf{Lax 函数}
\begin{align}
G(y,x,t)=\int_0^y u_0(z)\,\mathrm{d}z+tg\left(\frac{x-y}{t}\right),\label{align:ch11-11-4-8}
\end{align}
定义于 $(x,t)\in(-\infty,\infty)\times(0,\infty)$ 和 $y\in(-\infty,\infty)$。

\begin{theorem}\label{theorem:ch11-11-4-3}
对固定 $(x,t)\in(-\infty,\infty)\times(0,\infty)$, Lax 函数 $G(y,x,t)$ 在点 $\bar y\in(-\infty,\infty)$ 处取极小当且仅当连接点 $(x,t)$ 与 $(\bar y,0)$ 的直线段是与初值为 $u_0$ 的可容许解 $u$ 关联的广义特征线, 且在 $(0,t)$ 上无激波。
\end{theorem}

\begin{proof}
在 $(-\infty,\infty)$ 中固定 $y$ 与 $\bar y$, 在顶点为 $(x,t),(y,0),(\bar y,0)$ 的三角形上积分\eqref{align:ch11-11-1-1}并应用 Green 公式, 得
\begin{align}
&\int_0^{\bar y}u_0(z)\,\mathrm{d}z+\int_0^t\left[\frac{x-\bar y}{t}u\left(\bar y+\tau\frac{x-\bar y}{t}\pm,\tau\right)-f\left(u\left(\bar y+\tau\frac{x-\bar y}{t}\pm,\tau\right)\right)\right]\,\mathrm{d}\tau\nonumber\\
&\quad=\int_0^y u_0(z)\,\mathrm{d}z+\int_0^t\left[\frac{x-y}{t}u\left(y+\tau\frac{x-y}{t}\pm,\tau\right)-f\left(u\left(y+\tau\frac{x-y}{t}\pm,\tau\right)\right)\right]\,\mathrm{d}\tau.\label{align:ch11-11-4-9}
\end{align}
由\eqref{align:ch11-11-4-7}与\eqref{align:ch11-11-4-8}, \eqref{align:ch11-11-4-9}左端小于等于 $G(\bar y,x,t)$, 等号成立当且仅当在 $(0,t)$ 上几乎处处 $f'(u(\bar y+\tau\frac{x-\bar y}{t}\pm,\tau))=\frac{x-\bar y}{t}$, 即连接 $(x,t)$ 与 $(\bar y,0)$ 的直线段是无激波特征线。类似地, \eqref{align:ch11-11-4-9}右端小于等于 $G(y,x,t)$, 等号成立当且仅当连接 $(x,t)$ 与 $(y,0)$ 的直线段是无激波特征线。于是假设连接 $(x,t)$ 与 $(\bar y,0)$ 的直线段是无激波特征线, 由\eqref{align:ch11-11-4-9}推出对任意 $y\in(-\infty,\infty)$, $G(\bar y,x,t)\leqslant G(y,x,t)$。

反之, 假设对所有 $y\in(-\infty,\infty)$, $G(\bar y,x,t)\leqslant G(y,x,t)$。特别地, 取 $y$ 使 $(y,0)$ 是从 $(x,t)$ 出发的极小后向特征线在 $x$ 轴上的截距。如上所示, $y$ 是 $G(\cdot,x,t)$ 的极小点, 故 $G(y,x,t)=G(\bar y,x,t)$。此外, \eqref{align:ch11-11-4-9}右端等于 $G(y,x,t)$, 故左端也等于 $G(y,x,t)$。如上所释, 这蕴含连接 $(x,t)$ 与 $(\bar y,0)$ 的直线段是无激波特征线。证毕。
\end{proof}

上述命题可用于从初值 $u_0$ 确定可容许解 $u$: 对固定 $(x,t)\in(-\infty,\infty)\times(0,\infty)$, 设 $y_-$ 表示 $G(\cdot,x,t)$ 在 $(-\infty,\infty)$ 上的最小极小点, $y_+$ 表示最大极小点。则
\begin{align}
u(x\pm,t)=[f']^{-1}\left(\frac{x-y_\pm}{t}\right).\label{align:ch11-11-4-10}
\end{align}
由\cref{theorem:ch11-11-3-2}、\cref{theorem:ch11-11-3-3}与\cref{theorem:ch11-11-3-4}, 我们断定 $(x,t)$ 是 $u$ 的连续点当且仅当 $y_-=y_+$; 是 $u$ 的跳跃间断点当且仅当 $y_-<y_+$ 且 $y_-,y_+$ 是 $G(\cdot,x,t)$ 仅有的极小点; 是 $u$ 的不规则点当且仅当 $y_-<y_+$ 且在区间 $(y_-,y_+)$ 内还存在 $G(\cdot,x,t)$ 的其他极小点。可以基于上述可容许解构造发展真正非线性标量守恒律 Cauchy 问题的整套理论, 以替代广义特征线方法。然而应注意, 广义特征线方法提供了更大的灵活性, 因为它适用于定义在 $\mathbb{R}^2$ 任意开子集上的解, 而不必是整个上半平面。

变量代换 $u=\partial_x v$ 将守恒律\eqref{align:ch11-11-1-1}化为\textbf{Hamilton-Jacobi 方程}
\begin{align}
\partial_t v(x,t)+f(\partial_x v(x,t))=0.\label{align:ch11-11-4-11}
\end{align}
在此框架下, $u$ 是\eqref{align:ch11-11-1-1}的可容许弱解当且仅当 $v$ 是\eqref{align:ch11-11-4-11}的\textbf{粘性解}; 见\cref{chapter:ch11-一维标量守恒律}所引文献。事实上,\cref{theorem:ch11-11-4-2}与\cref{theorem:ch11-11-4-3}反映的是 Hamilton-Jacobi 方程解的性质, 而非双曲守恒律的性质, 因为它们可自然推广到前者的多空间维版本, 却不能推广到后者。


\end{document}
